\documentclass[11pt]{article}

\usepackage[margin=1in]{geometry}
\usepackage{amsmath,amssymb,mathtools,empheq}
\usepackage{booktabs}
\usepackage{microtype}
\usepackage[round]{natbib}
\usepackage[colorlinks=true,allcolors=blue]{hyperref}

\newcommand{\dd}{\mathrm{d}}
\newcommand{\Dt}{\mathcal{D}_t}
\newcommand{\pos}[1]{\left\langle #1\right\rangle_{+}}
\newcommand{\rmi}{\mathrm{i}}
\newcommand{\rms}{\mathrm{s}}
\newcommand{\rmb}{\mathrm{b}}
\newcommand{\KIc}{K_{\mathrm{Ic}}}

\title{A van der Veen Hydrofracture Criterion\\as a Damage-Evolution Law}
\author{Technical note}
\date{\today}

\begin{document}
\maketitle

\begin{abstract}
The van der Veen crevasse models give a fracture-propagation or crack-arrest
criterion, not a rate law.  Here the fractured fraction of an ice column is
written as \(D=D_{\rms}+D_{\rmb}\), where \(D_{\rms}\) and \(D_{\rmb}\) are
the fractions penetrated by surface and basal crevasses.  The criterion is
converted into an irreversible overstress law.  The central result is the
explicit source term
\[
 \Dt D
 =\pos{1-D}^{p}
 \left[
   \frac{1}{\tau_{\rms}}\pos{\Psi_{\rms}}^{m_{\rms}}
  +\frac{1}{\tau_{\rmb}}\pos{\Psi_{\rmb}}^{m_{\rmb}}
 \right],
\]
with either stress-intensity-factor drivers (the literal van der Veen LEFM
model) or inexpensive crack-tip-stress drivers (the common reduced
crevasse-depth model).  The new timescales and rate exponents are constitutive
parameters; they are not determined by the original static criterion.
\end{abstract}

\section{State variables and conventions}

Let \(H>0\) be the local ice thickness, \(d_{\rms}\) the downward penetration
of surface crevasses, and \(d_{\rmb}\) the upward penetration of basal
crevasses.  Define
\begin{equation}
 D_{\rms}=\frac{d_{\rms}}{H},\qquad
 D_{\rmb}=\frac{d_{\rmb}}{H},\qquad
 D=D_{\rms}+D_{\rmb},\qquad 0\leq D\leq1.
 \label{eq:damage-definition}
\end{equation}
Thus \(D\) is a geometric, column-integrated damage variable: it is the
fraction of the thickness intersected by surface or basal fractures.  The
constraint \(D\leq1\) assumes that the two fractured intervals do not overlap.
Full-thickness connection, and hence the local calving/failure condition, is
\(D=1\).

The material derivative following horizontal ice motion \(\boldsymbol u\) is
\begin{equation}
 \Dt(\,\cdot\,)=\frac{\partial(\,\cdot\,)}{\partial t}
                 +\boldsymbol u\boldsymbol\cdot\nabla(\,\cdot\,),
 \label{eq:material-derivative}
\end{equation}
and \(\pos{x}=\max(x,0)\).  Damage is irreversible here; healing can be
added as a separate, explicitly justified process if desired.

For a vertical fracture whose normal is the horizontal direction \(n\), let
\(R_n\) denote the tensile resistive stress.  In a three-dimensional shallow
stress model it should include the crevasse-parallel deviatoric stress; for
example, \(R_n=2\tau_{nn}+\tau_{tt}\), rather than using \(\tau_{nn}\) alone
\citep{reynolds2025}.  Tension is positive.

\section{The original criteria}

\subsection{Literal van der Veen LEFM criterion}

The surface and bottom-crevasse models of \citet{vanderveen1998surface} and
\citet{vanderveen1998bottom} use linear elastic fracture mechanics (LEFM).
For either crack \(j\in\{\rms,\rmb\}\), propagation occurs when
\begin{equation}
 K_j>\KIc,
 \label{eq:lefm-criterion}
\end{equation}
and an advancing crack arrests where \(K_j=\KIc\).  Here \(K_j\) is the net
mode-I stress intensity factor and \(\KIc\) is the fracture toughness of ice.

For completeness, write \(a_{\rms}=H D_{\rms}\), let \(y\) measure depth
downward from the upper surface, and let \(\delta_{\rms}\) be the dry length
above the water surface in a surface crack.  A convenient statement of the
van der Veen surface SIF is
\begin{align}
 K_{\rms}={}&F(\lambda_{\rms})R_n\sqrt{\pi a_{\rms}}
 -\frac{2\rho_{\rmi}g}{\sqrt{\pi a_{\rms}}}
   \int_0^{a_{\rms}} y\,
      G\!\left(\lambda_{\rms},\frac{y}{a_{\rms}}\right)\dd y \notag\\
 &+\frac{2\rho_{\rm m}g}{\sqrt{\pi a_{\rms}}}
   \int_0^{a_{\rms}}\pos{y-\delta_{\rms}}
      G\!\left(\lambda_{\rms},\frac{y}{a_{\rms}}\right)\dd y,
 \qquad \lambda_{\rms}=\frac{a_{\rms}}{H}.
 \label{eq:surface-sif}
\end{align}
The three terms are the contributions of resistive stress, ice overburden,
and surface meltwater.  The water-column height is
\(h_{\rms}=\pos{a_{\rms}-\delta_{\rms}}\), and \(\rho_{\rm m}\) is the
meltwater density.

For a basal crack, let \(y\) measure upward from the bed, let \(p_{\rmb}^0\)
be basal water pressure at the crack mouth, and let \(\rho_{\rmb}\) be the
water density within the crack.  The corresponding weight-function form is
\begin{align}
 K_{\rmb}={}&\frac{2}{\sqrt{\pi a_{\rmb}}}
 \int_0^{a_{\rmb}}
 \left[
 R_n-\rho_{\rmi}g(H-y)
 +\pos{p_{\rmb}^0-\rho_{\rmb}gy}
 \right]
 G\!\left(\lambda_{\rmb},\frac{y}{a_{\rmb}}\right)\dd y,
 \notag\\[-0.25em]
 &\hspace{30em}\lambda_{\rmb}=\frac{a_{\rmb}}{H}.
 \label{eq:basal-sif}
\end{align}
The finite-slab functions \(F\) and \(G\) are listed in
Appendix~\ref{sec:weights}.  Equations~\eqref{eq:surface-sif} and
\eqref{eq:basal-sif} make all hydrostatic assumptions explicit and are
directly usable in a quadrature routine.

LEFM requires an existing flaw.  Numerically, initialize
\(D_j\geq D_{j0}>0\), or evaluate \(K_j\) using
\(a_j=H\max(D_j,D_{j0})\).  The seed size \(D_{j0}\) is a model parameter,
not a prediction of LEFM.

\subsection{Reduced crack-tip-stress criterion}

A widely used, cheaper closure places the crack tip where the net opening
stress vanishes.  It is often grouped with ``van der Veen-type'' or
crevasse-depth hydrofracture laws, but it is not the same approximation as
the LEFM criterion above.  Define
\begin{equation}
 \Sigma=\frac{R_n}{\rho_{\rmi}gH},\qquad
 W_{\rms}=\frac{h_{\rms}}{H},\qquad
 r_{\rms}=\frac{\rho_{\rm m}}{\rho_{\rmi}},\qquad
 r_{\rmb}=\frac{\rho_{\rmb}}{\rho_{\rmi}},
 \label{eq:nondimensional-groups}
\end{equation}
and the basal height above water-pressure support
\begin{equation}
 H_{\rm ab}=H-\frac{p_{\rmb}^0}{\rho_{\rmi}g},
 \qquad A_{\rmb}=\frac{H_{\rm ab}}{H}.
 \label{eq:hab}
\end{equation}
For marine-connected basal water, \(p_{\rmb}^0=\rho_{\rmb}gw\), where \(w\)
is water depth at the bed.  The reduced basal expression below assumes that
water reaches the basal crack tip.  The dimensionless crack-tip opening
stresses are then
\begin{align}
 \Psi_{\rms}^{\rm ZS}
   &=\Sigma+r_{\rms}W_{\rms}-D_{\rms},
 \label{eq:surface-zs-driver}\\
 \Psi_{\rmb}^{\rm ZS}
   &=\Sigma-A_{\rmb}-(r_{\rmb}-1)D_{\rmb}.
 \label{eq:basal-zs-driver}
\end{align}
The label ZS denotes ``zero stress.''  Positive values open the crack; zero
values give the usual static depths
\begin{align}
 D_{\rms}^{\star}
   &=\Sigma+r_{\rms}W_{\rms},
 \label{eq:surface-target}\\
 D_{\rmb}^{\star}
   &=\frac{1}{r_{\rmb}-1}\pos{\Sigma-A_{\rmb}}
    =\frac{\rho_{\rmi}}{\rho_{\rmb}-\rho_{\rmi}}
      \pos{\frac{R_n}{\rho_{\rmi}gH}-\frac{H_{\rm ab}}{H}}.
 \label{eq:basal-target}
\end{align}
Equation~\eqref{eq:basal-target} assumes \(\rho_{\rmb}>\rho_{\rmi}\).
These are the familiar surface and basal crevasse-depth formulas used in
large-scale calving parameterizations \citep{nick2010,reynolds2025}.

\section{Damage-evolution law}

\subsection{Generic kinetic wrapper}

Define a dimensionless fracture driver for each crack family:
\begin{equation}
 \Psi_j^{\rm LEFM}=\frac{K_j}{\KIc}-1
 \quad\text{(literal van der Veen)},
 \qquad\text{or}\qquad
 \Psi_j=\Psi_j^{\rm ZS}
 \quad\text{(reduced model)}.
 \label{eq:driver-options}
\end{equation}
The proposed evolution equations are
\begin{empheq}[box=\fbox]{align}
 \Dt D_{\rms}
 &=\frac{\pos{1-D}^{p}}{\tau_{\rms}}
   \pos{\Psi_{\rms}}^{m_{\rms}},
 \label{eq:surface-evolution}\\
 \Dt D_{\rmb}
 &=\frac{\pos{1-D}^{p}}{\tau_{\rmb}}
   \pos{\Psi_{\rmb}}^{m_{\rmb}},
 \label{eq:basal-evolution}\\
 \Dt D
 &=\pos{1-D}^{p}
 \left[
   \frac{1}{\tau_{\rms}}\pos{\Psi_{\rms}}^{m_{\rms}}
  +\frac{1}{\tau_{\rmb}}\pos{\Psi_{\rmb}}^{m_{\rmb}}
 \right].
 \label{eq:total-evolution}
\end{empheq}
Here \(\tau_j>0\) is a propagation/relaxation timescale,
\(m_j\geq1\) is a rate-sensitivity exponent, and \(p>0\) controls smooth
saturation as the remaining ligament vanishes.  The minimal choice is
\begin{equation}
 m_{\rms}=m_{\rmb}=p=1.
 \label{eq:minimal-parameters}
\end{equation}

Equation~\eqref{eq:total-evolution} is the requested right-hand side for the
damage material derivative.  It has four useful properties:
\begin{enumerate}
 \item it is dimensionally correct because each \(\Psi_j\) is dimensionless
       and \(\tau_j\) has units of time;
 \item the positive part enforces irreversibility and reproduces the original
       propagation threshold;
 \item static crack-arrest states are unchanged: growth stops when the driver
       first reaches zero;
 \item the saturation factor makes the interval \(0\leq D\leq1\) invariant
       for the continuous equations.
\end{enumerate}
The original van der Veen work does not determine \(\tau_j\), \(m_j\), or
\(p\).  They must be calibrated, resolved as fast brittle timescales, or taken
to an instantaneous limit.  For \(p>0\), \(D=1\) is approached
asymptotically, so a numerical model should declare calving at
\(D\geq1-\epsilon_{\rm calv}\).  If finite-time arrival at exactly \(D=1\)
is required, omit the saturation factor for \(D<1\) and project the update
onto the admissible set after each step.

\subsection{Fully explicit reduced law}

Substituting equations~\eqref{eq:surface-zs-driver} and
\eqref{eq:basal-zs-driver} gives an algebraic right-hand side with no SIF
quadrature:
\begin{empheq}[box=\fbox]{align}
 \Dt D
 =\pos{1-D}^{p}\Bigg[&
 \frac{1}{\tau_{\rms}}
 \pos{\Sigma+r_{\rms}W_{\rms}-D_{\rms}}^{m_{\rms}}
 \notag\\[-0.15em]
 &+\frac{1}{\tau_{\rmb}}
 \pos{\Sigma-A_{\rmb}-(r_{\rmb}-1)D_{\rmb}}^{m_{\rmb}}
 \Bigg].
 \label{eq:explicit-reduced-law}
\end{empheq}
This is the most convenient form for a depth-integrated ice-flow model.  The
two component equations, rather than the total equation alone, must be
advanced because the right-hand side depends separately on \(D_{\rms}\) and
\(D_{\rmb}\).

For constant water-column height, the surface term relaxes toward
equation~\eqref{eq:surface-target}.  Hydrology can instead make
\(W_{\rms}\) a function of damage.  In particular, a fully water-filled
surface crack has \(W_{\rms}=D_{\rms}\), so
\begin{equation}
 \Psi_{\rms}^{\rm ZS}
 =\Sigma+(r_{\rms}-1)D_{\rms}.
 \label{eq:full-water-feedback}
\end{equation}
Because freshwater is denser than ice, \(r_{\rms}>1\); crack deepening then
increases the driver and produces the expected hydrofracture runaway.  A
water-volume or water-supply equation is therefore needed if
\(h_{\rms}\) is not prescribed.

\section{Coupling and implementation choices}

\paragraph{Stress--damage feedback.}
The \(R_n\) in the driver should be the trial resistive stress consistent with
the ice-flow solver.  If a fixed membrane force is carried only by the intact
ligament, a simple load-sharing closure is
\begin{equation}
 R_n^{\rm lig}=\frac{N_n}{H(1-D)},
 \label{eq:ligament-stress}
\end{equation}
which adds destabilizing feedback.  Do not apply
equation~\eqref{eq:ligament-stress} a second time if the constitutive or flow
solver already accounts for the loss of load-bearing area.

\paragraph{Thickness evolution.}
Equations~\eqref{eq:surface-evolution}--\eqref{eq:total-evolution} evolve the
fraction \(D_j\) as the material state.  If a code instead evolves dimensional
crack depths, the exact kinematic conversion is
\begin{equation}
 \Dt d_j=H\,\Dt D_j+D_j\,\Dt H.
 \label{eq:dimensional-depth-rate}
\end{equation}
One should decide whether changing thickness stretches the interpreted crack
depth or whether \(D_j\) is retained as a history variable; mixing the two
interpretations introduces a spurious damage source or sink.

\paragraph{Time stepping.}
At each material point: (i) compute trial stress and water state; (ii) evaluate
the two drivers; (iii) advance both component equations; (iv) enforce
non-negativity and cap \(D_{\rms}+D_{\rmb}\leq1\); and (v) apply the calving
threshold.  Explicit stepping should resolve the smaller of
\(\tau_{\rms}\) and \(\tau_{\rmb}\), or use substepping/implicit integration.

\paragraph{Geometry warning.}
The van der Veen weight functions describe a finite rectangular slab with a
single edge crack.  They are not universal weight functions; grounded basal
boundary conditions can materially change the SIF, especially for deep
cracks \citep{jimenez2018}.  For a grounded-glacier application, retain the
kinetic wrapper in equations~\eqref{eq:surface-evolution}--
\eqref{eq:total-evolution} but replace \(K_j\) by one computed with a
geometry-appropriate weight function or numerical fracture solution.

\section{Recommended minimal model}

For a first implementation in a large-scale model, use
equation~\eqref{eq:explicit-reduced-law} with
\(m_{\rms}=m_{\rmb}=p=1\), advect \(D_{\rms}\) and \(D_{\rmb}\) separately,
and calve at \(D\geq1-\epsilon_{\rm calv}\).  This gives the compact system
\begin{align}
 \Dt D_{\rms}
 &=\frac{1-D}{\tau_{\rms}}
   \pos{\Sigma+r_{\rms}W_{\rms}-D_{\rms}},
 \notag\\
 \Dt D_{\rmb}
 &=\frac{1-D}{\tau_{\rmb}}
   \pos{\Sigma-A_{\rmb}-(r_{\rmb}-1)D_{\rmb}},
 \qquad D=D_{\rms}+D_{\rmb}.
 \label{eq:recommended-model}
\end{align}
When literal adherence to van der Veen LEFM is required, keep exactly the same
evolution equations and replace the two bracketed drivers with
\(\pos{K_j/\KIc-1}\), using equations~\eqref{eq:surface-sif} and
\eqref{eq:basal-sif}.

\appendix
\section{van der Veen finite-slab weight functions}
\label{sec:weights}

The functions reproduced by \citet{jimenez2018} are
\begin{equation}
 F(\lambda)=1.12-0.23\lambda+10.55\lambda^2
             -21.72\lambda^3+30.39\lambda^4,
 \label{eq:F-weight}
\end{equation}
and
\begin{align}
 G(\lambda,\gamma)={}&
 \frac{3.52(1-\gamma)}{(1-\lambda)^{3/2}}
 -\frac{4.35-5.28\gamma}{(1-\lambda)^{1/2}}
 \notag\\
 &+\left[
   \frac{1.30-0.30\gamma^{3/2}}{(1-\gamma^2)^{1/2}}
   +0.83-1.76\gamma
 \right]\left[1-(1-\gamma)\lambda\right],
 \label{eq:G-weight}
\end{align}
where \(\lambda=a/H\) and \(\gamma=y/a\).  The crack-tip singularity in
\(G\) is integrable.  In computation, use endpoint-aware quadrature and stop
the LEFM calculation short of \(\lambda=1\), where the slab formula is
singular and the column has already met the calving criterion.

\bibliographystyle{plainnat}
\bibliography{references}

\end{document}
